\section*{Capítulo \ref{ch:b3:statistical-thermo-applied} --- Termodinámica estadística aplicada}

\begin{solution}{exo:b3:statistical-thermo-applied:1}
$\frac32R + R + (3 \times 3 - 5)R = 6.5R = \qty{54.0}{J.K^{-1}.mol^{-1}}$. A temperatura ambiente, la traslación
y la rotación son clásicas, las dos tensiones están congeladas y la flexión doblemente
degenerada, la vibración más baja, está parcialmente excitada.
\end{solution}

\begin{solution}{exo:b3:statistical-thermo-applied:2}
$x = 1$: $\eu/(\eu - 1)^2 = 0.921$. $x = 10$: $100\eu^{10}/(\eu^{10} - 1)^2 \approx 100\eu^{-10} = \num{4.5e-3}$.
\end{solution}

\begin{solution}{exo:b3:statistical-thermo-applied:3}
\qty{20}{K}: $x_{\mathrm{para}} = 1/(1 + 3 \times 3\eu^{-2 \times 85.35/20}) = 1/(1 + 9 \times \num{1.96e-4}) = 0.998$.
\qty{77}{K}: suma par $1 + 5\eu^{-6.651} = 1.0065$; suma impar $3(3\eu^{-2.217} + 7\eu^{-13.30}) = 0.9806$;
$x_{\mathrm{para}} = 1.0065/1.9871 = 0.51$.
\end{solution}

\begin{solution}{exo:b3:statistical-thermo-applied:4}
Para $I = 1$ hay $3 \times 2 = 6$ estados de espín simétricos y $3 \times 1 = 3$ antisimétricos;
los deuterones son bosones, de modo que la función de onda total es simétrica: los estados de espín
simétricos van con $J$ par (orto, peso 6) y los antisimétricos con $J$ impar (para, peso
3). El ortodeuterio, que contiene $J = 0$, es el estable a 0 K; a temperatura ambiente,
orto:para = 2:1.
\end{solution}

\begin{solution}{exo:b3:statistical-thermo-applied:5}
$\Delta_rE_0 = 2 \times 1883.7 - 2170.3 - 1542.3 = \qty{54.8}{cm^{-1}}$; $\eu^{-54.8/207.2} = 0.768$, y
$4 \times 0.768 = 3.07$. El valor completo, 3.26, es un 6\,\% mayor: los factores de traslación, rotación y
vibración solo se cancelan exactamente en el límite clásico, y la
rotación de estas moléculas ligeras sigue parcialmente cuantizada a \qty{298}{K}.
\end{solution}

\begin{solution}{exo:b3:statistical-thermo-applied:6}
\ce{N2O}: dos orientaciones, $R\ln2 = \qty{5.76}{J.K^{-1}.mol^{-1}}$. \ce{CH3D}: cuatro, $R\ln4 =
\qty{11.53}{J.K^{-1}.mol^{-1}}$ (límites superiores, alcanzados si las disposiciones son aleatorias).
\end{solution}

\begin{solution}{exo:b3:statistical-thermo-applied:7}
$x = 797.6/300 = 2.659$: función de Einstein $x^2\eu^{-x}/(1 - \eu^{-x})^2 = 0.572$; $C_V/R = 3.072$,
$C_p = 4.072R = \qty{33.86}{J.K^{-1}.mol^{-1}}$, un 0.4\,\% por debajo de la tabla (anarmonicidad).
\end{solution}

\begin{solution}{exo:b3:statistical-thermo-applied:8}
$x = 1.029$: $x^2\eu^{-x}/(1 - \eu^{-x})^2 = 0.916$, es decir, \qty{7.62}{J.K^{-1}.mol^{-1}}, el 92\,\% de $R$: la
vibración del yodo es casi clásica a temperatura ambiente.
\end{solution}

\begin{solution}{exo:b3:statistical-thermo-applied:9}
B es más rica en \ce{^{18}O}. $\alpha_{\mathrm{A/B}} = (1 - 0.0100)/(1 + 0.0020) = 0.98802$.
\end{solution}

\begin{solution}{exo:b3:statistical-thermo-applied:10}
$\Delta\mathrm{ZPE} = \frac12(1 - 1/\sqrt2)(2000 - 3000) = \qty{-146}{cm^{-1}}$; $K = \eu^{146.4/207.2} = 2.0$.
El deuterio va preferentemente a X, el enlace más rígido, donde su ahorro de \omterm{def:b3:quantum-model-systems:oscillator}{energía de
punto cero} es mayor.
\end{solution}

\begin{solution}{exo:b3:statistical-thermo-applied:11}
$\Delta_rH^\circ = R\ln10\,(-1.766 + 3.392)/(1/900 - 1/1100) = 19.145 \times 1.626/\num{2.020e-4} =
\qty{154}{kJ/mol}$. Los dos átomos llevan $2 \times \frac52RT$ de entalpía; la molécula, $\frac52RT + RT$
(rotación) más su energía de vibración $R\theta_{\mathrm V}/(\eu^{\theta_{\mathrm V}/T} - 1)$, algo menor que $RT$:
la diferencia, de unos \qty{5}{kJ/mol} a \qty{1000}{K}, se suma a $D_0$.
\end{solution}

\begin{solution}{exo:b3:statistical-thermo-applied:12}
Niveles 0 ($g = 1$) y $6hcB$ ($g = 5$): $q = 1 + 5\eu^{-x}$, $\langle\varepsilon\rangle/kT = 5x\eu^{-x}/q$, $\langle\varepsilon^2\rangle/(kT)^2 =
5x^2\eu^{-x}/q$, y $C/R = \langle\varepsilon^2\rangle/(kT)^2 - (\langle\varepsilon\rangle/kT)^2 = 5x^2\eu^{-x}/(1 + 5\eu^{-x})^2$. A
\qty{100}{K}, $x = 5.121$: $C/R = 0.783/1.061 = 0.738$.
\end{solution}

\begin{solution}{pb:b3:statistical-thermo-applied:1}
\textbf{1.} $2 \times 2 = 4$: tres simétricos (el triplete) y uno antisimétrico (el singlete).
\textbf{2.} Los protones son fermiones, de modo que la función de onda total cambia de signo con el
intercambio; la función de rotación da $(-1)^J$: $J$ par va con el singlete
antisimétrico (para), y $J$ impar con el triplete simétrico (orto).
\textbf{3.} Niveles pares 1, niveles impares 3.
\textbf{4.} A \qty{300}{K} ($T = 3.5\,\theta_{\mathrm R}$), las sumas de rotación par e impar son casi iguales;
ponderadas por 1 y 3 dan 1:3 (fracción de para 0.2507).
\textbf{5.} $F(1) = 2 \times 59.322 - 4 \times 0.0471 = \qty{118.46}{cm^{-1}} = \qty{1.417}{kJ/mol}$.
\textbf{6.} $1/(1 + 9\eu^{-170.4/20.37}) = 1/(1 + 9 \times \num{2.3e-4}) = 0.998$.
\textbf{7.} Con solo $J = 0$ y 1, la mitad de para significa $9\eu^{-2\theta_{\mathrm R}/T} = 1$, $T = 2\theta_{\mathrm R}/\ln9 =
0.91\,\theta_{\mathrm R} = \qty{78}{K}$: la competición es entre el peso 9 de $J = 1$ y su
factor de Boltzmann.
\textbf{8.} 0.75. Sin catalizador, la conversión necesita días; la licuefacción,
horas.
\textbf{9.} $(0.75 - 0.002) \times 1.417 = \qty{1.060}{kJ/mol}$.
\textbf{10.} $1.060/\num{2.016e-3} = \qty{526}{kJ/kg}$.
\textbf{11.} $0.905/\num{2.016e-3} = \qty{449}{kJ/kg}$.
\textbf{12.} El calor de conversión es 1.17 veces la entalpía de vaporización.
\textbf{13.} Se produce dentro del líquido; no entra a través de las paredes.
\textbf{14.} $x_{\mathrm o}(t) = x_{\mathrm{eq}} + (0.75 - x_{\mathrm{eq}})\eu^{-kt}$, con $x_{\mathrm{eq}} = 0.002$.
\textbf{15.} $0.002 + 0.748\eu^{-0.274} = 0.571$.
\textbf{16.} $(0.75 - 0.571) \times 1.417 = \qty{0.254}{kJ/mol}$.
\textbf{17.} $0.254/0.905 = 0.28$: el 28\,\% del depósito en un día.
\textbf{18.} Tras \qty{168}{h}, $x_{\mathrm o} = 0.112$, calor \qty{0.904}{kJ/mol}, igual a la entalpía de
vaporización: en este modelo, el depósito queda vacío. La estimación exagera la
pérdida porque las moléculas evaporadas se llevan consigo su \omterm{def:b3:statistical-thermo-applied:spin-isomers}{ortohidrógeno} sin
convertir; la pérdida real es menor, pero sigue siendo grande.
\textbf{19.} $t_{1/2} = \ln2/0.0114 = \qty{61}{h}$.
\textbf{20.} En el licuefactor, sobre lechos de catalizador a temperaturas sucesivas, de modo que el
refrigerador elimine el calor de conversión antes de almacenar el líquido.
\textbf{21.} Sí: la fracción de para en el equilibrio a \qty{300}{K} es 0.2507.
\textbf{22.} En el hidrógeno normal, las dos formas no pueden convertirse una en otra, de modo que la
rotación de cada una se congela en su propia escalera de niveles; en el hidrógeno en equilibrio, parte
del calor aportado convierte para en orto.
\textbf{23.} Alrededor de 3.57 hacia \qty{50}{K}: además de excitar la rotación, el calor convierte
moléculas de $J = 0$ en $J = 1$, \qty{118}{cm^{-1}} más arriba, una reacción de entalpía
positiva.
\textbf{24.} 1.50: la rotación está congelada en ambas formas (para en $J = 0$, orto en $J = 1$).
\textbf{25.} \textbf{Calor de conversión / entalpía de vaporización $\approx 1.2$ para el hidrógeno
normal}: su conversión por sí sola podría evaporar todo el depósito, y por eso el
hidrógeno se convierte en para antes de almacenarlo.
\end{solution}
